\section{Exact kernels of positive moment pairings}
\label{sec:exact-kernels}
The annihilation of a subspace and the realization of that subspace as
an entire kernel are different questions. The distinction is governed
by multiplication of functions, before a prior is chosen. We describe
it by a linear quotient and a determinantal test, and then construct
full-support witnesses arbitrarily close to any admissible prior.

Let $X$ be a nonempty compact metric space, and let $E,F\subset C(X;\R)$
have dimensions $d,k$. No ordering of these dimensions is required in
this section. Measures $\mu$ below are Borel probabilities. For $U\subset F$ put
\begin{equation}\label{eq:kernel-product-spaces}
 \begin{gathered}
 W_U=\spanop\{eu:e\in E,\ u\in U\},\qquad
 V=\spanop(\{1\}\cup EF),\\
 \mathcal P_U=\{\mu:\operatorname{supp}\mu=X,\ \mu W_U=0\}.
 \end{gathered}
\end{equation}
Here $EF$ denotes the set of pointwise products, not a tensor product.
The condition $\mu W_U=0$ means that every member of $W_U$ has integral
zero. It is equivalent to $U\subset\ker H_\mu$, where
$H_\mu f=(\mu(e_a f))_{a=1}^d$ in a fixed basis of $E$.

Suppose for the moment that $\mathcal P_U$ is nonempty. Then
$1\notin W_U$. Choose $v_1,\ldots,v_s\in V$ so that
\[
 [1],[v_1],\ldots,[v_s]\quad\hbox{is a basis of }V/W_U,
 \qquad s=\dim V-\dim W_U-1.
\]
Choose representatives $\bar f_1,\ldots,\bar f_q$ of a basis of $F/U$,
where $q=k-\dim U$. There are unique coefficients with
\begin{equation}\label{eq:kernel-pencil}
 e_a\bar f_i\equiv c_{ai0}+\sum_{\nu=1}^s c_{ai\nu}v_\nu
       \pmod {W_U},\qquad
 \mathscr H_U(z)=C_0+\sum_{\nu=1}^s z_\nu C_\nu.
\end{equation}
Thus $\mathscr H_U$ is a $d$ by $q$ affine matrix pencil determined by
the multiplication table modulo $W_U$. Define its algebraic rank by
\begin{equation}\label{eq:kernel-generic-rank}
 r_U=\rank_{\R(z_1,\ldots,z_s)}\mathscr H_U.
\end{equation}
For $s=0$ this is the rank of the constant matrix. For $q=0$ it is zero.
A change of bases makes an invertible affine change of moment coordinates
and invertible row or column operations, so $r_U$ is intrinsic.

\begin{lemma}[Local coordinates on a constrained prior slice]
\label{lem:kernel-moment-chart}
If $\mu_*\in\mathcal P_U$, the set
\[
 \mathcal O_U=\{(\mu v_1,\ldots,\mu v_s):\mu\in\mathcal P_U\}
\]
is a nonempty open convex subset of $\R^s$. At every $\mu_*$ it admits
a local parametrization by bounded continuous density perturbations
that preserve all the equations $\mu W_U=0$. More precisely, there are
$h_1,\ldots,h_s\in C(X)$ and a positive definite matrix $B$ such that
\begin{equation}\label{eq:kernel-tilt-chart}
 d\mu_t=(1+\textstyle\sum_\nu t_\nu h_\nu)d\mu_*,\qquad
 \mu_t v=\mu_*v+Bt,\qquad \mu_t W_U=0,
\end{equation}
and these are full-support probabilities whenever
$\|\sum_\nu t_\nu h_\nu\|_\infty<1$.
\end{lemma}
\begin{proof}
Full support makes the $L^2(\mu_*)$ inner product positive definite on
any finite-dimensional space of continuous functions: a nonzero
continuous function has nonzero magnitude on a nonempty open set.
Since $\mu_*W_U=0$, that space is orthogonal to the constants. Let
$\Pi_U$ be its $L^2(\mu_*)$ orthogonal projection and set
\[
 h_\nu=v_\nu-\mu_*v_\nu-\Pi_U v_\nu.
\]
These are continuous, have mean zero, and are orthogonal to $W_U$.
They are linearly independent, since a relation among them would put
the same linear combination of the $v_\nu$ in $\R1+W_U$.
Consequently $B_{\alpha\nu}=\mu_*(h_\alpha h_\nu)$ is positive
definite. Orthogonality gives
$\mu_*(v_\alpha h_\nu)=B_{\alpha\nu}$, proving all identities in
\eqref{eq:kernel-tilt-chart}. The strict uniform bound on the tilt
gives a strictly positive density, normalization follows from the
zero means, and full support is retained. Since $B$ is invertible,
the image contains a neighborhood of $\mu_*v$. Repeating this argument
at every point proves openness. Convex mixtures of full-support
probabilities prove convexity. When $s=0$, the moment space is the
one-point space $\R^0$ and the assertion uses the empty family of tilts.
\end{proof}

\begin{theorem}[Exact-kernel criterion and full-support realization]
\label{thm:exact-kernel}
A subspace $U\subset F$ is the kernel of $H_\mu$ for some full-support
probability if and only if the following two conditions hold:
\begin{equation}\label{eq:exact-kernel-criterion}
 W_U\cap C(X;\R)_+=\{0\},\qquad r_U=k-\dim U.
\end{equation}
The first condition is equivalent to $\mathcal P_U\ne\varnothing$
by Lemma~\ref{lem:dimension-free-feasibility}, including $W_U=0$.
It permits the definition of the pencil in \eqref{eq:kernel-pencil}. The second condition says that at
least one maximal column minor of that pencil is a nonzero polynomial.
When $q=0$, the empty minor is one.

More generally, subject to the first condition,
\begin{equation}\label{eq:kernel-slice-max-rank}
 \max_{\mu\in\mathcal P_U}\rank H_\mu=r_U.
\end{equation}
The priors attaining this maximum form a relatively weakly open dense
subset of $\mathcal P_U$. For every $\mu_*\in\mathcal P_U$ and
$0<\varepsilon<1$, such a prior can be chosen with continuous density
$1+h$ relative to $\mu_*$ and $\|h\|_\infty<\varepsilon$.
In the chart of Lemma~\ref{lem:kernel-moment-chart}, a tensor grid with
$r_U+1$ points per coordinate inside any sufficiently small box contains
a witness. This is a finite algebraic test of at most $(r_U+1)^s$
points, not an assertion that unspecified real moments are computable.

When \eqref{eq:exact-kernel-criterion} holds, for any prescribed
full-support probability $\mu_0$ an exact witness can also be chosen as
$\delta\mu_0+(1-\delta)\nu$, where $0<\delta<1$ and $\nu$ has at
most $\dim V$ atoms. No common positive value of $\delta$ is asserted.
\end{theorem}
\begin{proof}
Lemma~\ref{lem:dimension-free-feasibility}, applied to $W=W_U$,
proves the equivalence of the first condition and
$\mathcal P_U\ne\varnothing$ without comparing $d$ and $k$.
This includes $W_U=0$. It also gives $1\notin W_U$:
a probability annihilating $W_U$ cannot annihilate $1$.

For $\mu\in\mathcal P_U$, the matrix induced by $H_\mu$ on $F/U$
is exactly $\mathscr H_U(\mu v)$. All minors of order $r_U+1$
vanish identically, and some minor of order $r_U$ is a nonzero
polynomial, unless $r_U=0$, when every induced matrix is zero.
A nonzero real polynomial cannot vanish on a nonempty open set. By
Lemma~\ref{lem:kernel-moment-chart} the nonzero minor therefore has a
nonzero value at an admissible full-support moment. This proves
\eqref{eq:kernel-slice-max-rank}. Rank--nullity now proves
\eqref{eq:exact-kernel-criterion}. Necessity follows from the same
matrix identity and the fact that there are only $q$ columns.

For the stronger local assertion fix $\mu_*$. In its chart, the chosen
minor becomes a nonzero polynomial in $t$, since $B$ is invertible.
Its degree is at most $r_U$. Choose a box so small that every density
perturbation in it has uniform norm less than $\varepsilon$. A
polynomial of degree at most $r_U$ in each variable that vanishes on a
product of $r_U+1$ distinct values per variable is identically zero:
apply the univariate root bound to the last variable and induct on the
number of variables. Hence the grid contains a nonzero value of the
minor. The perturbation retains full support and every annihilation
equation. Arbitrarily small boxes prove density, even by uniform density
perturbations. Maximal rank is relatively weakly open because the
entries are integrals of continuous functions and a nonzero minor stays
nonzero nearby. If $r_U=0$, all priors are witnesses and the grid may
consist of its center. The same convention handles $s=0$.

Finally, take an exact full-support witness already constructed and
record all its moments in a basis of $V$ beginning with $1$. The
nonconstant moment body has affine dimension $\dim V-1$: an affine
relation would be a linear dependence in that basis. Apply
Lemma~\ref{lem:v19-moment-image} to this vector of moments and the
prescribed $\mu_0$. Its finite-atom representation has at most $\dim V$
atoms and preserves every product in $EF$. The whole pairing matrix,
and therefore its exact kernel, is unchanged.
\end{proof}

The polynomial in this criterion is tested as an identity on the full
affine quotient, not by an unverified search on a constrained set of
priors. Positivity and the annihilation equations have already been
resolved by the first condition and the local chart. The argument uses
finite-dimensional separation, orthogonal projection, and elementary
determinantal geometry; it is not a new generic-matrix rank theorem.
Its content here is the exact compatibility with full support, prescribed
annihilation, and arbitrarily small positive prior perturbations.

\subsection{Forced directions and a positive block classification}
A useful preliminary test is the multiplication closure
\begin{equation}\label{eq:kernel-closure}
 \operatorname{cl}_E(U)=\{f\in F: ef\in W_U\text{ for every }e\in E\}.
\end{equation}
It records directions forced to vanish by the annihilation equations
alone, whereas the pencil in \eqref{eq:kernel-pencil} retains all
remaining determinantal restrictions.

\begin{proposition}[The forced kernel]
\label{prop:forced-kernel}
If $\mathcal P_U\ne\varnothing$, then
\begin{equation}\label{eq:forced-kernel-intersection}
 \bigcap_{\mu\in\mathcal P_U}\ker H_\mu=\operatorname{cl}_E(U).
\end{equation}
Moreover $W_{\operatorname{cl}_E(U)}=W_U$, and this closure is
idempotent. In particular, exact realization of $U$ requires
$\operatorname{cl}_E(U)=U$. The exact sufficient condition is the
maximal-minor test in Theorem~\ref{thm:exact-kernel}.
\end{proposition}
\begin{proof}
Every member of the right side has all its products in $W_U$, so is
annihilated by every admissible prior. Conversely, let $f$ belong to
every kernel and fix $e\in E$. Put $g=ef$ and choose
$\mu_*\in\mathcal P_U$. Then $\mu_*g=0$. Orthogonally project $g$
onto $W_U$ in $L^2(\mu_*)$ and put $h=g-\Pi_U g$. It has mean
zero, is orthogonal to $W_U$, and is bounded and continuous. For all
sufficiently small real $t$, the measure $(1+th)\mu_*$ belongs to
$\mathcal P_U$. Its integral of $g$ is $t\mu_*(h^2)$, which must
vanish. Thus $h=0$ in $L^2(\mu_*)$, hence identically by continuity
and full support. Therefore $g\in W_U$. Since this holds for every
$e$, the claimed equality follows. The inclusions
$U\subset\operatorname{cl}_E(U)$ and
$E\operatorname{cl}_E(U)\subset W_U$ give equality of the product
spaces. Applying the definition again proves idempotence.
\end{proof}

The following family makes the closure and the exact criterion explicit.
It includes the failure of an unqualified kernel-realization claim, but
also classifies the kernels that are realized in the same positive
experiment.

\begin{proposition}[Exact balanced-block kernels]
\label{prop:block-kernels}
Let $n\ge m+1$, let $J=\{1,\ldots,n\}$, and put
\[
 X=J\times\{-1,1\},\quad E=\spanop\{e_j:j\in J\},\quad
 F=\R1\oplus\sigma\mathbb R_{<m}[z].
\]
Here $e_j$ is the indicator of the $j$th two-point block and a polynomial
$p$ represents the function $(j,\sigma)\mapsto p(j)$. Let
$0\ne S\subset\mathbb R_{<m}[z]$, $U=\sigma S$, and
\[
 Z_S=\{j\in J:p(j)=0\text{ for every }p\in S\}.
\]
Then $\mathcal P_U\ne\varnothing$,
\begin{equation}\label{eq:block-kernel-classification}
 \operatorname{cl}_E(U)
   =\sigma\{p\in\mathbb R_{<m}[z]:p|_{Z_S}=0\},\qquad
 r_U=1+|Z_S|.
\end{equation}
Thus $U$ is exactly realizable if and only if
\begin{equation}\label{eq:block-root-condition}
 S=\left(\prod_{j\in Z_S}(z-j)\right)
                  \mathbb R_{<m-|Z_S|}[z].
\end{equation}
All these assertions have a positive affine acquisition and a fixed
finite physical query realization.
\end{proposition}
\begin{proof}
Write $s_j=\mu(j,+)+\mu(j,-)>0$ and
$d_j=\mu(j,+)-\mu(j,-)$, so $|d_j|<s_j$ and $\sum_j s_j=1$.
The equations $\mu E U=0$ are $d_jp(j)=0$ for every $p\in S$.
They force $d_j=0$ on $J\setminus Z_S$ and impose no condition on
$d_j$ on $Z_S$. The uniform probability is an admissible witness.
Since a nonzero polynomial of degree less than $m$ has at most $m-1$
roots, $|Z_S|\le m-1$ and $J\setminus Z_S$ is nonempty.

In the basis $1,\sigma,\sigma z,\ldots,\sigma z^{m-1}$ the rows of
the pairing are
\[
 (s_j,d_j,d_jj,\ldots,d_jj^{m-1}).
\]
Any kernel element $c+\sigma p$ has $c=0$, by a row outside $Z_S$.
For fixed $d$, its remaining constraint is $p(j)=0$ at those $j\in Z_S$
where $d_j\ne0$. Distinct-point polynomial evaluation has rank equal
to the number of such points, because that number is less than $m$.
It follows that the maximal pairing rank is $1+|Z_S|$. It is attained
explicitly by $s_j=1/n$, $d_j=1/(2n)$ on $Z_S$, and $d_j=0$
elsewhere. All atom masses are positive. The kernel of this witness is
exactly the first space in \eqref{eq:block-kernel-classification}; every
admissible kernel contains that space. Proposition~\ref{prop:forced-kernel}
identifies it with the closure. Rank--nullity proves the exact-realization
criterion, and the factor theorem for polynomials gives
\eqref{eq:block-root-condition}.

For the physical realization use $n-1$ acquisition features
\[
 f_i(j,\sigma)=\frac{e_i(j,\sigma)}{2(n-1)},\qquad
 L_y(u,j,\sigma)=\tfrac12(1+y u^tf(j,\sigma)),
 \quad u\in[-1,1]^{n-1}.
\]
They satisfy $L_y\ge1/4$. At $u=0$, the history product and its command
derivatives span $E$, since $e_n=1-\sum_{i<n}e_i$. The query
probabilities
\[
 g_\ell(j,\sigma)=\tfrac12+\tfrac14\sigma(j/n)^\ell,
 \qquad 0\le\ell<m,
\]
belong to $[1/4,3/4]$, and their span with $1$ is $F$. They can each
be executed as a binary query with the stated success probability.
The detector, acquisition cube, and query menu are fixed while the
prior varies. This proves the physical assertion.
\end{proof}

For $n=4$, $m=2$, and $S=\spanop\{1\}$, the set $Z_S$ is empty.
Thus $U=\spanop\{\sigma\}$ satisfies the annihilation criterion but
its forced closure is $\spanop\{\sigma,j\sigma\}$; the pairing rank
on the entire constrained slice is one. The product-evaluation vectors
are $\pm\mathbf e_j$, so their convex hull contains zero in its
interior. There is nevertheless no exact kernel $U$. In contrast, for
each $j_0\in J$, the line $\spanop\{(j-j_0)\sigma\}$ satisfies
\eqref{eq:block-root-condition} and is an exact kernel of the explicit
full-support witness above. The relative-interior condition and the
exact-kernel condition therefore make different, verifiable predictions
inside the same experiment.

\subsection{Return to the acquired prediction state}
\begin{corollary}[Exact constrained prediction ranks]
\label{cor:kernel-history-rank}
Let $P:O\to C(X;\R)$ be a positive history-product map of class $C^2$,
where $u_*$ is an interior point of the command neighborhood $O$, and let
$\Phi=\spanop\{1,\phi_1,\ldots,\phi_b\}$ be a finite continuous
future test space. Apply the preceding construction with
\[
 F=\spanop\{P(u_*,\cdot)\}+\operatorname{im}DP(u_*),\qquad E=\Phi.
\]
For every $U\subset F$ with $\mathcal P_U\ne\varnothing$, the maximal
normalized prediction rank at $u_*$ among these priors is $r_U-1$.
It is attained by the arbitrarily small bounded-density perturbations
in Theorem~\ref{thm:exact-kernel}. If the exact-kernel condition holds,
the attained rank is $\dim F-\dim U-1$.

For each maximizing prior with $r_U>1$, commands with a density
bounded below near $u_*$
and the specified report word give an unconditional subprobability
minorization on an $(r_U-1)$-dimensional projection of predictions.
\end{corollary}
\begin{proof}
Let $L:F\to E^*$ be the moment pairing and let
$Y=LP=(Z,Zp)^t$, where $Z=\mu P>0$. Normalization has kernel
$\R Y$ on the image of $L$. Moreover
\[
 L(F)=\R Y+L(\operatorname{im}DP(u_*)).
\]
Hence the actual command derivative of the normalized map has rank
$\rank L-1$; the evidence column need not itself be a command
derivative. The positive evidence also ensures $\rank L\ge1$.
Theorem~\ref{thm:exact-kernel} gives the maximal rank and its local
attainment, and rank--nullity gives the exact-kernel specialization.

For a maximizing prior choose $r_U-1$ independent output rows of the
command derivative and complete them by orthonormal kernel coordinates.
The resulting square map has an invertible derivative. A sufficiently
small product box in its image has inverse contained in the command
neighborhood, and its inverse Jacobian determinant is bounded below.
The joint density of the commands and the specified report word is the
command density times $\mu P(u)$, which has a positive lower bound on
a smaller neighborhood. Change variables and integrate over the entire
kernel-coordinate box. The resulting positive lower density on the
output box proves the asserted subprobability minorization. It is a
statement about actual histories, not mass assigned to a section.
\end{proof}

For the balanced-block witnesses, normalization yields the prediction
rank $|Z_S|$. Their full one-step resolution law is then supplied by
Theorem~\ref{thm:v18-affine-classification}, using the covariance of the
fixed acquisition and query functions. In the general multi-step problem,
the present kernel criterion supplies local attainment information; the
whole-image cover and stable causal update remain separate requirements.
In the monomial experiment these are precisely the other two parts of
Theorems~\ref{thm:intrinsic-checkpoint} and
\ref{thm:intrinsic-streaming}.
